% Representative primary-source nodes; line spacing is not proportional to years.
\begin{tikzpicture}[
  x=1mm,y=1mm,font=\figurefont\fontsize{8.5}{11}\selectfont,
  event/.style={draw=BookInk,fill=white,line width=.8pt,
    minimum width=37mm,minimum height=14mm,align=center,inner sep=1.5mm}
]
  \foreach \y/\name/\idea in {
    48/试错学习/行为因结果而改变,
    24/动态规划/长期价值可以递归,
    0/时序差分预测/相邻预测提供误差
  } {
    \node[anchor=east,font=\figurefont\bfseries\fontsize{9}{11}\selectfont]
      at (25,\y) {\name};
    \draw[BookMuted,line width=.8pt] (30,\y) -- (124,\y);
    \node[anchor=west,align=left,text width=28mm] at (131,\y) {\idea};
  }
  \node[event] at (52,48) {Thorndike（1911）\\效果律与试错};
  \node[event] at (103,48) {Barto等（1983）\\学习控制结构};
  \node[event] at (52,24) {Bellman（1957）\\动态规划};
  \node[event] at (103,24) {Howard（1960）\\MDP与策略迭代};
  \node[event] at (52,0) {Samuel（1959）\\经验与局面评价};
  \node[event] at (103,0) {Sutton（1988）\\TD预测分析};
\end{tikzpicture}
